% در این فایل، دستورها و تنظیمات مورد نیاز، آورده شده است.
%-------------------------------------------------------------------------------------------------------------------
% دستوراتی که پوشه پیش‌فرض زیرفایل‌های tex را مشخص می‌کند.
%\makeatletter
%\def\input@path{{./tex/}}
%\makeatother
% در ورژن جدید زی‌پرشین برای تایپ متن‌های ریاضی، این سه بسته، حتما باید فراخوانی شود
\usepackage{amsthm,amssymb,amsmath}
% بسته‌ای برای تنطیم حاشیه‌های بالا، پایین، چپ و راست صفحه
\usepackage[a4paper, top=40mm, bottom=40mm, outer=25mm, inner=35mm]{geometry}
% بسته‌‌ای برای ظاهر شدن شکل‌ها و تعیین آدرس تصاویر
\usepackage[final]{graphicx}
\graphicspath{{./img/}}



% بسته‌های مورد نیاز برای نوشتن کدها، رنگ‌آمیزی آنها و تعیین پوشه کدها
%%man ezafe kardam
\usepackage{subfig}

\usepackage{diagbox}
\usepackage{multirow}
\usepackage{rotating}
\usepackage[final]{listings}
\usepackage[usenames,dvipsnames,svgnames,table]{xcolor}
\lstset{inputpath=./code/}
\usepackage[notransparent]{svg}
\usepackage[final]{pdfpages}
\usepackage{tabularray}
\usepackage{booktabs, makecell, multirow, tabularx, longtable}
% بسته‌ای برای رسم کادر
\usepackage{framed}
% بسته‌‌ای برای چاپ شدن خودکار تعداد صفحات در صفحه «معرفی پایان‌نامه»
\usepackage{lastpage}
% بسته‌ لازم برای: ۱. تغییر شماره‌گذاری صفحات پیوست. ۲. تصحیح باگ آدرس وب حاوی '%' در مراجع
\usepackage{etoolbox}

%%%%%%%%%%%%%%%%%%%%%%%%%%%%%%%%%%%%
%%% دستورات وابسته به استیل مراجع:
%% اگر از استیل‌های natbib (plainnat-fa، asa-fa، chicago-fa) استفاده می‌کنید، خط زیر را فعال و بعدی‌اش را غیرفعال کنید.
%\usepackage{natbib}
%\newcommand{\citelatin}[1]{\cite{#1}\LTRfootnote{\citeauthor*{#1}}}
%\newcommand{\citeplatin}[1]{\citep{#1}\LTRfootnote{\citeauthor*{#1}}}


%% اگر از سایر استیل‌ها استفاده می‌کنید، خط بالا را غیرفعال و خط‌های زیر را فعال کنید.
\let\citep\cite
\let\citelatin\cite
\let\citeplatin\cite
%%%%%%%%%%%%
%% بررسی حالت پیش نویس
\usepackage{ifdraft}
\ifdraft
{%
	% بسته‌ ایجاد لینک‌های رنگی با امکان جهش
	\usepackage[unicode=true,pagebackref=true,
colorlinks,linkcolor=blue,citecolor=blue,final]{hyperref}
	%\usepackage{todonotes}
	\usepackage[firstpage]{draftwatermark}
	\SetWatermarkText{\ \ \ \rl{پیش‌نویس}}
	\SetWatermarkScale{1.2}
}
{
	\usepackage[pagebackref=false]{hyperref}
	%\usepackage[disable]{todonotes} % final without TODOs
}

\usepackage[obeyDraft]{todonotes}
\setlength{\marginparwidth}{2cm}

%%%%%%%%%%%%
%%% تصحیح باگ: اگر در مراجع، آدرس وب حاوی '%' بوده و pagebackref فعال باشد، دستورات زیر باید بیایند:
%% برای استیل‌های natbib مثل plainnat-fa، asa-fa، chicago-fa
\makeatletter
\let\ORIG@BR@@lbibitem\BR@@lbibitem
\apptocmd\ORIG@BR@@lbibitem{\endgroup}{}{}
\def\BR@@lbibitem{\begingroup\catcode`\%=12 \ORIG@BR@@lbibitem}
\makeatother
%% برای سایر استیل‌ها
\makeatletter
\let\ORIG@BR@@bibitem\BR@@bibitem
\apptocmd\ORIG@BR@@bibitem{\endgroup}{}{}
\def\BR@@bibitem{\begingroup\catcode`\%=12 \ORIG@BR@@bibitem}
\makeatother
%%%%%%%%%%%%%%%%%%%%%%%%%%%%%%%%%%%%

% بسته‌ لازم برای تنظیم سربرگ‌ها
\usepackage{fancyhdr}
\setlength{\headheight}{16.1pt}
% Give Persian paragraphs a small amount of flexibility so long mixed-script
% terms do not protrude into the margins.
\setlength{\emergencystretch}{3em}
%\usepackage{enumitem}
\usepackage{setspace}
% بسته‌های لازم برای نوشتن الگوریتم
\usepackage{algorithm}
\usepackage{algorithmic}
% بسته‌های لازم برای رسم بهتر جداول
\usepackage{tabulary}
\usepackage{tabularx}
\usepackage{rotating}
% بسته‌های لازم برای رسم تنظیم بهتر شکل‌ها و زیرشکل‌ها
\usepackage[export]{adjustbox}
\usepackage{subfig}
\usepackage[subfigure]{tocloft}
% کپشن‌ها و محتوای جدول‌ها اندکی کوچک‌تر از متن اصلی باشند.
\captionsetup{font=footnotesize,labelfont=bf}
\AtBeginEnvironment{tabular}{\footnotesize}
\AtBeginEnvironment{tabularx}{\footnotesize}
\AtBeginEnvironment{longtable}{\footnotesize}
% بسته‌ای برای رسم نمودارها و نیز صفحه مالکیت اثر
\usepackage{tikz}
\usetikzlibrary{arrows.meta,positioning,shapes.geometric}
\usetikzlibrary{calc,fit,backgrounds,shapes.misc,shapes.multipart}
\usetikzlibrary{decorations.pathreplacing,decorations.markings,patterns}
\usepackage{pgfplots}
\pgfplotsset{compat=1.17}
\usepgfplotslibrary{groupplots}

%%% تنظیم پارامترهای جای‌گذاری شناورها.
%%% با افزایش تعداد شکل‌ها و جدول‌ها، تنظیمات پیش‌فرض لاتک باعث انباشت
%%% شناورهای جای‌گذاری‌نشده و در نهایت خطای «Too many unprocessed floats» می‌شود.
\renewcommand{\topfraction}{0.85}
\renewcommand{\bottomfraction}{0.70}
\renewcommand{\textfraction}{0.12}
\renewcommand{\floatpagefraction}{0.70}
\setcounter{topnumber}{3}
\setcounter{bottomnumber}{2}
\setcounter{totalnumber}{5}

%%%%%%%%%%%%%%%%%%%%%%%%%%%%%%%%%%%%%%%%%%%%%%%%%%%%%%%%%%%%%%%%%%%%%%%%%%%%%%
%%% پالت رنگ و سبک یکنواخت شکل‌ها و نمودارهای پایان‌نامه.
%%% هدف: تمام شکل‌ها با یک زبان بصری واحد خوانده شوند؛ هر مدل همیشه
%%% با یک رنگ ثابت نمایش داده می‌شود تا خواننده نیازی به بازخوانی
%%% راهنمای هر نمودار نداشته باشد.
%%%%%%%%%%%%%%%%%%%%%%%%%%%%%%%%%%%%%%%%%%%%%%%%%%%%%%%%%%%%%%%%%%%%%%%%%%%%%%
\definecolor{thProposed}{HTML}{1F6FB2} % مدل پیشنهادی (با ماتریس کالیبراسیون دوربین)
\definecolor{thControl}{HTML}{C0392B} % مدل کنترل (بدون ماتریس کالیبراسیون دوربین)
\definecolor{thBase}{HTML}{78838C} % مدل پایه منتشرشده
\definecolor{thIdeal}{HTML}{2E9E5B} % مرجع/رفتار ایده‌آل هندسی
\definecolor{thAccent}{HTML}{E08A1E} % تاکید و حاشیه‌نویسی
\definecolor{thInk}{HTML}{2B2B2B} % رنگ خط و متن شکل‌ها
\definecolor{thPanel}{HTML}{EAF1F8}
\definecolor{thPanelAlt}{HTML}{FDF0E6}
\definecolor{thPanelGreen}{HTML}{E7F4EC}
\definecolor{thPanelGray}{HTML}{F0F1F2}
\definecolor{thPanelRed}{HTML}{FBEAE8}

\pgfplotsset{
 thesisplot/.style={
  width=0.86\textwidth,
  height=6.4cm,
  grid=major,
  grid style={gray!25, very thin},
  axis line style={gray!55},
  tick align=outside,
  label style={font=\small},
  tick label style={font=\small},
  title style={font=\small},
  legend cell align=left,
  legend style={font=\small, draw=gray!45, fill=white},
  every axis plot/.append style={thick},
 },
}

\tikzset{
 thblock/.style={draw=thInk!60, thick, rounded corners=2pt, align=center,
         inner sep=5pt, font=\small},
 thdata/.style={draw=thInk!45, thick, rounded corners=9pt, align=center,
         inner sep=5pt, font=\small, fill=thPanelGray},
 thnote/.style={font=\scriptsize, align=center, text=thInk},
 tharrow/.style={-{Latex[length=2.2mm]}, thick, draw=thInk!75},
 thdashed/.style={-{Latex[length=2.2mm]}, thick, dashed, draw=thAccent},
}
% بسته‌ای برای ظاهر شدن «مراجع» و «نمایه» در فهرست مطالب
\usepackage[nottoc,notlof,notlot]{tocbibind}
% دستورات مربوط به ایجاد نمایه
\usepackage{makeidx}
\makeindex


% بسته‌ای برای افزودن تورفتگی به ابتدای اولین پاراگراف هر بخش
\usepackage{indentfirst}

% بسته زیر باگ ناشی از فراخوانی بسته‌های زیاد را برطرف می‌کند.
\usepackage{morewrites}

%%% سازگاری با نسخه‌های جدید array/LaTeX:
%%% وصله \@tabular در بسته bidi هنوز \UseMathForPositioningText را صدا می‌زند،
%%% ولی array.sty در توزیع‌های جدید آن را حذف کرده است. تعریف تهی زیر همان چیزی
%%% است که array نسخه ۲۰۲۴ ارایه می‌کرد؛ در توزیع‌های قدیمی‌تر بی‌اثر است.
\providecommand\UseMathForPositioningText{}
%%%%%%%%%%%%%%%%%%%%%%%%%%
%khodam
\RequirePackage{zref-perpage}
\zmakeperpage{footnote}

%%% درج هزارگان برای اعداد
%% Example: \num{52325} = 52,325
\usepackage{siunitx}
\sisetup{group-separator = {,}}


%%% بسته ایجاد واژه‌نامه با xindy
\usepackage[xindy,toc,acronym,nonumberlist=true]{glossaries}

% فراخوانی بسته زی‌پرشین (باید آخرین بسته باشد)
\usepackage[extrafootnotefeatures, displaymathdigits=persian]{xepersian}
\newcommand{\fafootnote}[1]{\footnote{#1}}
\newcommand{\enfootnote}[1]{\LTRfootnote{#1}}
% شماره پاورقی در هر دو جهت متن با رقم فارسی نمایش داده شود. جهت خود
% پاورقی را xepersian تعیین می‌کند: پاورقی فارسی راست‌به‌چپ و پاورقی انگلیسی
% چپ‌به‌راست باقی می‌ماند؛ این تعریف فقط قلم رقم را مستقل از جهت متن می‌کند.
\renewcommand{\thefootnote}{\persianfont\arabic{footnote}}
\renewcommand{\thempfootnote}{\persianfont\arabic{mpfootnote}}
% Persian decimal numbers need an LTR direction without switching to the
% Latin font, whose glyph set does not include Persian digits.
\newcommand{\fadecimal}[1]{\mbox{\beginL#1\endL}}
% نام‌های معیار برای مدل‌های اصلی پژوهش؛ ویژگی نسخه‌های خاص (انکودر،
% پروتکل، مجموعه‌داده یا دور آموزش) پس از این نام ثابت افزوده می‌شود.
\newcommand{\proposedmodel}{مدل پیشنهادی}
\newcommand{\controlmodel}{مدل کنترل}
\newcommand{\basemodel}{مدل پایه \lr{Depth Anything V2}}
\newcommand{\publishedbasemodel}{\basemodel{}، نسخه منتشرشده}
\newcommand{\calibrationmatrix}{\mbox{\persianfont\rl{ماتریس کالیبراسیون دوربین}}}
% برچسب‌های کوتاه ویژه جدول‌ها. جزییات پیکربندی (انکودر، پروتکل و
% داده آموزش) در ستون‌های مستقل یا عنوان جدول می‌آیند تا ستون «مدل» فشرده بماند.
\newcommand{\tableproposedmodel}{پیشنهادی}
\newcommand{\tablecontrolmodel}{کنترل}
\newcommand{\tablepublishedmodel}{پایه منتشرشده}
% \paragraphfootnotes

\makeatletter
% تعریف قلم فارسی و انگلیسی و مکان قلم‌ها
\if@irfonts
\settextfont[Path={./font/}, BoldFont={IRLotusICEE_Bold.ttf}, BoldItalicFont={IRLotusICEE_BoldIranic.ttf}, ItalicFont={IRLotusICEE_Iranic.ttf},Scale=1.2]{IRLotusICEE.ttf}
% LiberationSerif or FreeSerif as free equivalents of Times New Roman
\setlatintextfont[Path={./font/}, BoldFont={LiberationSerif-Bold.ttf}, BoldItalicFont={LiberationSerif-BoldItalic.ttf}, ItalicFont={LiberationSerif-Italic.ttf},Scale=1]{LiberationSerif-Regular.ttf}
% چنانچه می‌خواهید اعداد در فرمول‌ها، انگلیسی باشد، خط زیر را غیرفعال کنید
% و گزینه displaymathdigits=persian را از خط ۱۰۹ حذف کنید.
\setdigitfont[Path={./font/}, Scale=1.2]{IRLotusICEE.ttf}
% تعریف قلم‌های فارسی و انگلیسی اضافی برای استفاده در بعضی از قسمت‌های متن
\setiranicfont[Path={./font/}, Scale=1.3]{IRLotusICEE_Iranic.ttf}				% ایرانیک، خوابیده به چپ
\setmathsfdigitfont[Path={./font/}]{IRTitr.ttf}
%\defpersianfont\titlefont[Path={./font/}, Scale=1]{IRTitr.ttf}
\defpersianfont\titlefont[Path={./font/}, Scale=1]{IRLotusICEE.ttf}
% برای تعریف یک قلم خاص عنوان لاتین، خط بعد را فعال و ویرایش کنید و خط بعد از آن را غیرفعال کنید.
% \deflatinfont\latintitlefont[Scale=1]{LiberationSerif}
\font\latintitlefont=cmssbx10 scaled 2300 %cmssbx10 scaled 2300
\else
\settextfont{XB Niloofar}
\setlatintextfont{Junicode}
% چنانچه می‌خواهید اعداد در فرمول‌ها، انگلیسی باشد، خط زیر را غیرفعال کنید
% و گزینه displaymathdigits=persian را از خط ۱۰۹ حذف کنید.
\setdigitfont{XB Niloofar}
% تعریف قلم‌های فارسی و انگلیسی اضافی برای استفاده در بعضی از قسمت‌های متن
% \setmathsfdigitfont{XB Titre}
\defpersianfont\titlefont{XB Titre}
\deflatinfont\latintitlefont[Scale=1.1]{Junicode}
\fi
\makeatother

% برای استفاده از قلم نستعلیق خط بعد را فعال کنید.
\defpersianfont\nast[Path={./font/}, Scale=2]{IranNastaliq}
% فونت‌های جدید را در این جا وارد کنید.
%\defpersianfont\newfont[Path={./font/}, Scale=2]{newfont}

%%%%%%%%%%%%%%%%%%%%%%%%%%
% راستچین شدن todonotes
\presetkeys{todonotes}{align=right,textdirection=righttoleft}{}
\makeatletter
\providecommand\@dotsep{5}
\def\listtodoname{فهرست کارهای باقیمانده}
\def\listoftodos{\noindent{\Large\vspace{10mm}\listtodoname}\@starttoc{tdo}}
\renewcommand{\@todonotes@MissingFigureText}{شکل}
\renewcommand{\@todonotes@MissingFigureUp}{شکل}
\renewcommand{\@todonotes@MissingFigureDown}{جاافتاده}
\makeatother
% دستوری برای حذف کلمه «چکیده»
% \renewcommand{\abstractname}{}
% دستوری برای حذف کلمه «abstract»
%\renewcommand{\latinabstract}{}
% دستوری برای تغییر نام کلمه «اثبات» به «برهان»
\renewcommand\proofname{برهان}
% دستوری برای تغییر نام کلمه «کتاب‌نامه» به «مراجع»
\renewcommand{\bibname}{مراجع}
% دستوری برای تعریف واژه‌نامه انگلیسی به فارسی
\newcommand\persiangloss[2]{#1\dotfill\lr{#2}\\}
% دستوری برای تعریف واژه‌نامه فارسی به انگلیسی
\newcommand\englishgloss[2]{#2\dotfill\lr{#1}\\}
% تعریف دستور جدید «\پ» برای خلاصه‌نویسی جهت نوشتن عبارت «پروژه/پایان‌نامه/رساله»
\newcommand{\پ}{پروژه/پایان‌نامه/رساله }

%\newcommand\BackSlash{\char`\\}

%%%%%%%%%%%%%%%%%%%%%%%%%%
% \SepMark{-}

% تعریف و نحوه ظاهر شدن عنوان قضیه‌ها، تعریف‌ها، مثال‌ها و ...
\theoremstyle{definition}
\newtheorem{definition}{تعریف}[section]
\theoremstyle{theorem}
\newtheorem{theorem}[definition]{قضیه}
\newtheorem{lemma}[definition]{لم}
\newtheorem{proposition}[definition]{گزاره}
\newtheorem{corollary}[definition]{نتیجه}
\newtheorem{remark}[definition]{ملاحظه}
\theoremstyle{definition}
\newtheorem{example}[definition]{مثال}

%\renewcommand{\theequation}{\thechapter-\arabic{equation}}
%\def\bibname{مراجع}
\numberwithin{algorithm}{chapter}
\def\listalgorithmname{فهرست الگوریتم‌ها}
\def\listfigurename{فهرست تصاویر}
\def\listtablename{فهرست جداول}

% دستور های لازم برای تعریف ترجمه دستورات الگوریتم
\makeatletter
\renewcommand{\algorithmicrequire}{\if@RTL{ورودی:}\else Require:\fi}
\renewcommand{\algorithmicensure}{\if@RTL{خروجی:}\else Ensure:\fi}
\renewcommand{\algorithmicend}{\if@RTL{پایان}\else end\fi}
\renewcommand{\algorithmicif}{\if@RTL{اگر}\else if\fi}
\renewcommand{\algorithmicthen}{\if@RTL{آنگاه}\else then\fi}
\renewcommand{\algorithmicelse}{\if@RTL{وگرنه}\else else\fi}
\renewcommand{\algorithmicfor}{\if@RTL{برای}\else for\fi}
\renewcommand{\algorithmicforall}{\if@RTL{برای هر}\else for all\fi}
\renewcommand{\algorithmicdo}{\if@RTL{انجام بده}\else do\fi}
\renewcommand{\algorithmicwhile}{\if@RTL{تا زمانی که}\else while\fi}
\renewcommand{\algorithmicloop}{\if@RTL{تکرار کن}\else loop\fi}
\renewcommand{\algorithmicrepeat}{\if@RTL{تکرار کن}\else repeat\fi}
\renewcommand{\algorithmicuntil}{\if@RTL{تا زمانی که}\else until\fi}
\renewcommand{\algorithmicprint}{\if@RTL{چاپ کن}\else print\fi}
\renewcommand{\algorithmicreturn}{\if@RTL{بازگردان}\else return\fi}
\renewcommand{\algorithmicand}{\if@RTL{و}\else and\fi}
\renewcommand{\algorithmicor}{\if@RTL{و یا}\else or\fi} % TODO add better translate
\renewcommand{\algorithmicxor}{\if@RTL{یا}\else xor\fi} % TODO add better translate
\renewcommand{\algorithmicnot}{\if@RTL{نقیض}\else not\fi}
\renewcommand{\algorithmicto}{\if@RTL{تا}\else to\fi}
\renewcommand{\algorithmicinputs}{\if@RTL{ورودی‌ها}\else inputs\fi}
\renewcommand{\algorithmicoutputs}{\if@RTL{خروجی‌ها}\else outputs\fi}
\renewcommand{\algorithmicglobals}{\if@RTL{متغیرهای عمومی}\else globals\fi}
\renewcommand{\algorithmicbody}{\if@RTL{انجام بده}\else do\fi}
\renewcommand{\algorithmictrue}{\if@RTL{درست}\else true\fi}
\renewcommand{\algorithmicfalse}{\if@RTL{نادرست}\else false\fi}
\renewcommand{\algorithmicendif}{\algorithmicend شرط \algorithmicif}
\renewcommand{\algorithmicendfor}{\algorithmicend حلقه \algorithmicfor}
\renewcommand{\algorithmicendwhile}{\algorithmicend حلقه \algorithmicwhile}
\renewcommand{\algorithmicendloop}{\algorithmicend حلقه \algorithmicloop}
\renewcommand{\algorithmiccomment}[1]{\{{ #1}\}}
\makeatletter

%%%%%%%%%%%%%%%%%%%%%%%%%%%%
%%% دستورهایی برای سفارشی کردن سربرگ صفحات:
%\newcommand{\SetHeader}[1]{
% دستور زیر معادل با گزینه twoside است.
%\csname@twosidetrue\endcsname
\pagestyle{fancy}
%% دستورات زیر سبک صفحات fancy را تغییر می‌دهد:
% O=Odd, E=Even, L=Left, R=Right
% در هر دو حالت oneside و twoside، سربرگ بر اساس جای فیزیکی روی صفحه تنظیم می‌شود.
% سمت چپ: شماره و نام فصل؛ سمت راست: عنوان پایان‌نامه.
\fancyhead{}
\fancyhead[L]{\small\leftmark}
\fancyhead[R]{\parbox[b]{0.45\textwidth}{\raggedleft\footnotesize\@title}}
\fancyfoot{}
\fancyfoot[C]{\thepage}
\setlength{\headheight}{26pt}
%\fancyhead[ER]{\small\rightmark}
%\fancyhead[OR]{\footnotesize}
%\fancyhead[EL]{\footnotesize\rightmark}
\renewcommand{\headrulewidth}{0.75pt}
% شکل‌دهی شماره و عنوان فصل در سربرگ
\renewcommand{\chaptermark}[1]{\markboth{\@chapapp~\thechapter:\ #1}{}}
\makeatletter
%\renewcommand{\rightmark}[1]{\@title}
%\makeatother
%}
% در فایل هر فصل، پس از دستور \chapter، صفحه عنوان با
% «\thispagestyle{fancy}\clearpage» بسته می‌شود تا متن از صفحه بعد آغاز شود.
%%%%%%%%%%%%%%%%%%%%%%%%%%%%
%\def\MATtextbaseline{1.5}
%\renewcommand{\baselinestretch}{\MATtextbaseline}
\doublespacing
%%%%%%%%%%%%%%%%%%%%%%%%%%%%%
% دستوراتی برای اضافه کردن کلمه «فصل» در فهرست مطالب

\newlength\mylenprt
\newlength\mylenchp
\newlength\mylenapp

\renewcommand\cftpartpresnum{\partname~}
\renewcommand\cftchappresnum{\chaptername~}
\renewcommand\cftchapaftersnum{:}

\settowidth\mylenprt{\cftpartfont\cftpartpresnum\cftpartaftersnum}
\settowidth\mylenchp{\cftchapfont\cftchappresnum\cftchapaftersnum}
\settowidth\mylenapp{\cftchapfont\appendixname~\cftchapaftersnum}
\addtolength\mylenprt{\cftpartnumwidth}
\addtolength\mylenchp{\cftchapnumwidth}
\addtolength\mylenapp{\cftchapnumwidth}

\setlength\cftpartnumwidth{\mylenprt}
\setlength\cftchapnumwidth{\mylenchp}

\makeatletter
{\def\thebibliography#1{\chapter*{\refname\@mkboth
  {\uppercase{\refname}}{\uppercase{\refname}}}\list
  {[\arabic{enumi}]}{\settowidth\labelwidth{[#1]}
  \rightmargin\labelwidth
  \advance\rightmargin\labelsep
  \advance\rightmargin\bibindent
  \itemindent -\bibindent

  \listparindent \itemindent
  \parsep \z@
  \usecounter{enumi}}
  \def\newblock{}
  \sloppy
  \sfcode`\.=1000\relax}}

%اگر مایلید در شماره گذاری حرفی و ابجد به جای آ از الف استفاده شود دستورات زیر را فعال کنید.
%\def\@Abjad#1{%
% \ifcase#1\or الف\or ب\or ج\or د%
% \or هـ\or و\or ز\or ح\or ط%
% \or ی\or ک\or ل\or م\or ن%
% \or س\or ع\or ف\or ص%
% \or ق\or ر\or ش\or ت\or ث%
% \or خ\or ذ\or ض\or ظ\or غ%
% \else\@ctrerr\fi}
%
% \def\abj@num@i#1{%
% \ifcase#1\or الف\or ب\or ج\or د%
% \or هـ‍\or و\or ز\or ح\or ط\fi

% \ifnum#1=\z@\abjad@zero\fi}
%
% \def\@harfi#1{\ifcase#1\or الف\or ب\or پ\or ت\or ث\or

% ج\or چ\or ح\or خ\or د\or ذ\or ر\or ز\or ژ\or س\or ش\or ص\or ض\or ط\or ظ\or ع\or غ\or

% ف\or ق\or ک\or گ\or ل\or م\or ن\or و\or ه\or ی\else\@ctrerr\fi}

%
\makeatother

%%% امکان درج کد در سند
% در این قسمت رنگ، قلم و قالب‌بندی قسمت‌های مختلف یک کد تعیین می‌شود.
\lstdefinestyle{myStyle}{
	basicstyle=\ttfamily, % whole listing /w verbatim font
	keywordstyle=\color{blue}, % bold black keywords
	identifierstyle=, % nothing happens
	commentstyle=\color{LimeGreen}, % green comments
	stringstyle=\ttfamily\color{red}, % red typewriter font for strings
	showstringspaces=false % no special string spaces
	breaklines=true,
	breakatwhitespace=false,
	numbers=right, % line number formats
	numberstyle=\footnotesize\lr,
	numbersep=-10pt,
	frame=single,
	captionpos=b,
	captiondirection=RTL
}
\lstset{style=myStyle} % command to set default style
\def\lstlistingname{\rl{برنامه}}
\def\lstlistlistingname{\rl{فهرست برنامه‌ها}}


% for numbering subsubsections
\setcounter{secnumdepth}{3}
%to include subsubsections in the table of contents
\setcounter{tocdepth}{3}

\makeatletter
\renewcommand{\p@subfigure}{\thefigure.}
\makeatother
